\documentclass[a4paper]{article}
\usepackage{luatexja}
\usepackage{amsmath}
\usepackage{amssymb}
\usepackage{amsthm}
\usepackage{geometry}
\usepackage{hyperref}

\newcommand{\Jnorm}{J_{\mathrm{norm}}}
\newtheorem{lemma}{補題}
\newtheorem{proposition}[lemma]{命題}
\newtheorem{definition}[lemma]{定義}
\renewcommand{\proofname}{証明}

\geometry{margin=25mm}

\title{中性子の\(\beta\)崩壊はローター類の置換である：\\
クォーク電荷の台帳}

\author{https://github.com/hypernumbernet}
\date{\today}

\begin{document}
\maketitle

\begin{abstract}
中性子の\(\beta\)崩壊を、二重時空理論の二つの核子類の置換として読む。離散の同伴稿は、核を陽子類 \([R_{\rm p}]\) と中性子類 \([R_{\rm n}]\) の順序つき積として表し、各類にバリオン数 \(B=1\) と双対位相の巻き数 \(w=3\) を与える。強い力の同伴稿は、中性子を電磁ラベル \(e=0\) の場合とし、陽子をそのラベルが立つ場合とする。自由中性子の崩壊は、積の一つの因子を \([R_{\rm n}]\) から \([R_{\rm p}]\) へ置き換え、因子の個数と巻き数を保ったまま、その因子の電磁ラベルを \(+e\) へ上げる遷移である。総電荷を保つ放出は、電子殻の同伴稿がすでに用いる電荷 \(-e\) の電子である。電子のエネルギースペクトルが連続であることは、電気的に沈黙した第三の物体を実験室が要求するという事実であり、六生成子空間はその物体を置く場所をもつ。置換そのものは第三の物体を強制しない。両方の核子類が釣り合った軌跡 \(J=0\) に座れるからである。

クォーク電荷 \(+2/3\) と \(-1/3\) は、三つの加法的なラベルと二つの値だけで、二つの核子電荷 \(+1\) と \(0\) を再現する唯一の有理解である。その台帳は二つの類の整数電荷を三つのラベルへ割り付けた座標である。親理論の三つの空間軸は、同一の \(3+3\) 分割の三つの埋め込みであり、相異なる軸の射影子は可換でなく、単位へ和しない。したがって軸は、同時に立つ三つの色でも、三つの世代でもない。コンパクト \(\mathrm{SU}(2)\) は双対ローターである。色の個数因子は、この分割からは読めない。

自由な質量差 \(m_n-m_p\) は \(1293/1000\,\mathrm{MeV}\) と \(1294/1000\,\mathrm{MeV}\) のあいだに厳密にある。強い力の同伴稿が置く第一の節は \(14/19\,\mathrm{fm}\) と \(213/280\,\mathrm{fm}\) のあいだに厳密にあり、その半径でのクーロン・エネルギー \(ke^{2}/r\) は \(189/100\,\mathrm{MeV}\) より厳密に大きい。正の電気的自己エネルギーは荷電の類を重くする。自由な順序は逆であり、間隙はその自己エネルギーより小さい。類の順序は節での自己エネルギーではない。平均寿命は \(870\,\mathrm{s}\) と \(890\,\mathrm{s}\) のあいだにあり、\(\hbar/(m_n-m_p)c^{2}\) は \(6\times 10^{-22}\,\mathrm{s}\) より小さい。比は \(10^{24}\) を超える。置換は \(\gamma_s\) の節を横切る遷移ではない。

等スケール軌跡の上では、磁気平方は符号なし質量に等しく、電気平方もそれに等しい。そこにある正の質量は、電磁ラベルが立っていても立っていなくても、磁気をもつ。ラベルが零のとき、電荷とベクトルポテンシャルの積は双対角を動かさず、磁気の項は静止した物体に力を及ぼさない。一つの軸を保つ生成子は、平行な磁気ローターである。その軸上の双対位相 \(2\pi\) はローターを \(-1\) へ送り、双対位相 \(4\pi\) は単位元へ戻す。自由な間隙は、なお入力であり、その軌跡の上では磁気平方の超過である。

V--A 結合、寿命、スペクトルの形、\(g_A/g_V\)、フェルミ定数、ニュートリノ質量、世代、カビボ--小林--益川行列、構造関数、比 \(R\)、質量差の絶対値、および核磁子で書いたいずれの能率も導出しない。
\end{abstract}

\section{はじめに}
\label{sec:intro}

親理論は四つの相互作用を一つの幾何の二つのチャンネルとして読む~\cite{dst-main}。重力チャンネルは慣性質量に結びついた対 \((\phi,\theta)\) であり、電磁チャンネルは電荷に結びついた対 \((\alpha,\beta)\) である。強い力は、核密度における重力チャンネルの層状相として提案されている。色荷は、双対ローターの位相が担うねじれ荷の誤読として名づけられ、結合 \(\alpha_s\) は離散床 \(\varepsilon_N\) とは同定されない。電磁気と弱いチャンネルは同じ六生成子空間を占め、双対写像も双対ローターも V--A 結合を供給しない。

離散の同伴稿は、各核子を三層のねじれ星として埋め込む~\cite{dst-dio}。標準類 \([R_{\rm p}]\) と \([R_{\rm n}]\) は \(B=1\) と \(w=3\) を運び、質量数 \(A\) の核はそれらの代表元の順序積である。代数的安定は、その積が単数群に属すること、巻き数が三層トポロジーと両立する整数 \(\ell\) に対して \(w=3\ell\) であること、および \(|\Jnorm|\le 1\) であることを要求する。強い力の同伴稿は、同じ二類を核力の粒子として用いる~\cite{dst-strong}。中性子も両方のチャンネルを担う。電荷は零であるから電気の項はなく、重力の項が残る。陽子は電磁項も保つ。陽子と中性子は異なる類であるから、一つの間隔にある陽子--中性子対は、同種フェルミオンの席の排他だけでは禁じられない。

本稿は、その二類のあいだの遷移を中性子の\(\beta\)崩壊として読む。第\ref{sec:chart}節は、クォーク電荷が二つの整数電荷の台帳であることを示す。第\ref{sec:axes}節は、三つの軸をその台帳の三つの粒子から分ける。第\ref{sec:classes}節は、二つの類がすでに固定している量を集める。第\ref{sec:decay}節は置換の保存量である。第\ref{sec:gap}節は、自由な質量差が節でのクーロン・エネルギーでないことを示す。第\ref{sec:moment}節は、核子の電荷に沈黙した磁石であり、続いてクォークの台帳が能率に与える比と、測定された比との間隙である。第\ref{sec:open}節は、弱いチャンネルについて本稿が導出しないものを列挙する。

\section{クォーク電荷の台帳}
\label{sec:chart}

二つの核子の実験室電荷は、陽子の \(+e\) と中性子の \(0\) である。クォーク模型は、この二つを三つの加法的な電荷の和として書き、その電荷が二つの値 \(u\) と \(d\) だけをとると仮定する。陽子を二つと一つ、中性子を一つと二つに割り当てると、台帳は連立
\begin{equation}
\label{eq:chart}
2u+d=1,
\qquad
u+2d=0
\end{equation}
である。単位は \(e\) である。

\begin{lemma}
\label{lem:chart}
有理数の対 \((u,d)\) で~\eqref{eq:chart} を満たすものは \(u=2/3\)、\(d=-1/3\) だけである。
\end{lemma}

\begin{proof}
二式の差は \(u-d=1\) である。第二式から \(u=-2d\) であり、したがって \(-2d-d=1\)、すなわち \(d=-1/3\)、\(u=2/3\) である。
\end{proof}

この対は、三つのラベルの和が三であるあらゆる割り当てを整数にするわけではない。和が \(+1\) または \(0\) になるのは、割り当てが \((2,1)\) または \((1,2)\) であるときに限る。三つの \(u\) は \(+2\) であり、三つの \(d\) は \(-1\) である。台帳が自由粒子の電荷を整数に保つのは、核子をその二つの割り当てに制限するからである。中間子の台帳は、値とその反符号を一つずつ足し、やはり整数を与える。\(u\) と \(\bar d\) の和は \(+1\)、\(u\) と \(\bar u\) の和は \(0\) である。分数は、制限された割り当ての中で互いに消し合う。

補題が示すのは、分数電荷が二つの核子電荷から一意に逆算されるということである。三つの加法的な担体という仮定を置けば、実験室が測る整数は \(2/3\) と \(-1/3\) を強制する。仮定を置かなければ、強制されるのは整数のままである。二重時空理論の核子は一つのローター類であり、三つの担体の積ではない。第\ref{sec:axes}節が、三つの軸をその担体として読めない理由を述べる。台帳は、二つの類の電荷を三つのラベルで書いた座標として残る。ラベルを粒子として伝播させる関係、すなわち構造関数、ジェット、および比 \(R\) は、補題の外にある。

比 \(R\) の台帳値は、色の個数を別に置いたうえで \(R=N_c\sum_f q_f^2\) である。\(u\)、\(d\)、\(s\) と \(N_c=3\) では \(R=2\) である。個数を外すと和は \(2/3\) である。本稿はこの高原を再計算しない。個数因子を三つの軸から読むこともできない。次節の障害による。

\section{三つの軸}
\label{sec:axes}

親理論は、一つの空間軸 \(e_a\) に相対して、六つの生成子を可換な三つ組 \((B_a,E_{a+1},E_{a+2})\) と反可換な三つ組 \((E_a,B_{a+1},B_{a+2})\) とに分ける~\cite{dst-main}。可換な三つ組は \(\mathfrak{so}(2,1)\) であり、\(\mathfrak{su}(2)\) ではない。平方が \(-1\) である三つ組は磁気の三つ組だけであり、それらは四元数の表 \(IJ=K\) を満たす。この六空間のコンパクト \(\mathrm{SU}(2)\) は双対ローターである。

三つの空間軸は、それぞれそのような分割を与える。付随する射影子は冪等である。相異なる軸の射影子は可換でなく、三つの右射影子は単位へ和しない。親理論はそれらを、同一の \(3+3\) 幾何の代替的な埋め込みとして述べ、三つの世代ラベルとしては述べない。同じ文が色を閉じる。同時に立つ三つの色は、互いに可換で、一つの表現の中で単位を尽くすラベルである。軸の射影子は、その二つの性質をもたない。したがって三つの軸は三つのクォークではなく、三つの色でもない。

巻き数 \(w=3\) は、この障害の別の側にある。離散の同伴稿が核子に割り当てるのは、三つの粒子のバリオン数 \(1/3\) ではなく、一つの類の双対位相が三層で閉じることである。一つの軸を外へ送り出す遷移は、順序積の因子の個数を変えない。軸は因子ではない。漸近状態として離散理論が数えるのは類 \([R_{\rm p}]\) と \([R_{\rm n}]\) であり、軸ではない。台帳の分数が自由粒子として現れないのは、この意味で、閉じ込めるべき粒子が積の中に無いからである。これは閉じ込めポテンシャルの導出ではない。強い力の同伴稿が節の内側で述べる指数的な遮蔽は、核子に働く力の別の言明であり、本稿はその遮蔽を漸近自由性と同定しない。

世代についても同じ障害が先にある。親理論は、有限な離散ローター像を世代の個数とは読まない。ミュー粒子とタウ粒子は、核子の残りの二軸ではない。荷電パイオンがミュー粒子を電子より多く出す分岐は、世代の事実であり、本稿の置換からは決まらない。

\section{二つの核子類}
\label{sec:classes}

本稿が用いる類の性質は、同伴稿がすでに固定したものに限る。

\begin{enumerate}
\item どちらもバリオン数 \(B=1\) と巻き数 \(w=3\) をもつ~\cite{dst-dio}。一つの核子は三層で閉じ、巻き数 \(w=3\ell\) の整数 \(\ell\) は核子では \(\ell=1\) である。
\item 中性子の電磁ラベルは \(e=0\) であり、陽子の電磁ラベルは \(+e\) である~\cite{dst-strong}。電荷はチャンネルのラベルであり、不整合 \(J\) の符号ではない。
\item 核物質の読みでは、等スケール軌跡の上で \(J=0\) かつ \(M>0\) である~\cite{dst-main}。釣り合った種は、一方のラピディティ予算だけでは空にならない。
\item 電子は電磁チャンネルの粒子であり、その静止は二つの予算の一致が保つ~\cite{dst-shell}。水素のイオン化はその静止を支払わない。
\end{enumerate}

クラスの差としてまだ固定されていないのは、符号なし質量 \(M\) の差である。自由中性子が自由陽子より重いことは測定であり、次節では入力として用いる。軸ごとの質量は定義しない。台帳が \(m_d>m_u\) と書く超過は、本稿では類 \([R_{\rm n}]\) の超過であり、ラベル \(d\) の超過ではない。

順序積に第三の核子類は無い。ストレンジネスを変える崩壊は、本稿が導入しない類を要する。ハイペロンは本稿の外にある。

\section{類の置換としての\texorpdfstring{\(\beta\)}{β}崩壊}
\label{sec:decay}

\begin{definition}
\label{def:replacement}
核の順序積における\emph{類の置換}とは、一つの因子 \([R_{\rm n}]\) を一つの因子 \([R_{\rm p}]\) で置き換えること、またはその逆である。
\end{definition}

\begin{proposition}
\label{prop:conserved}
類の置換は、因子の個数 \(A\) を保ち、各因子の巻き数条件 \(w=3\) を保つ。中性子から陽子への置換は、その因子の電磁ラベルを \(0\) から \(+e\) へ変える。逆はラベルを \(+e\) から \(0\) へ戻す。
\end{proposition}

\begin{proof}
順序積の定義により、\(A\) は因子の個数である。置換は一つの因子を一つの因子で置き換えるので、個数は不変である。両方の類が \(w=3\) をもつので、閉じたねじれループの条件は、置換の前に成り立っていたならば置換の後にも成り立つ。電磁ラベルは強い力の同伴稿が二類に与えたものである。
\end{proof}

バリオン数を因子の個数として読むかぎり、置換はバリオン数を保つ。\(B=1/3\) のラベルを一つ変える必要はない。巻き数が保たれるので、崩壊は三層の一つをほどく遷移でもない。

総電荷が保存するならば、中性子から陽子への置換は電荷 \(-e\) を積の外へ出す。その電荷は、電子殻の同伴稿の電子がすでに担うラベルである。電子は核子の断片ではない。置換が核子の電磁ラベルを立てるとき、総和を保つために外へ出る電磁ローターである。

逆置換には二つの実験室的な実現がある。積が陽電子を出す実現と、すでに存在する電子を吸収する実現である。前者は \(\beta^{+}\) 崩壊であり、後者は電子捕獲である。どちらも命題~\ref{prop:conserved} の逆であり、電磁ラベルを沈黙へ戻す。分岐比は導出しない。

第三の物体は、置換の保存則からは出ない。電子の運動エネルギースペクトルが連続であることは、電気的に沈黙した粒子が電子とともに出ることを実験室が要求する。六生成子空間は、電荷ラベルをもたない双対セクターの励起を、その粒子の読みとして置く場所をもつ。強制は別である。

\begin{proposition}
\label{prop:no-third}
両方の核子類が \(J=0\) に座れるとき、置換が持ち出す不整合は零である。電荷をもつローターは \(J=0\)、\(M>0\) に座れる。したがって、釣り合った電荷の放出は、不整合を打ち消す第二のローターを要求しない。
\end{proposition}

\begin{proof}
親理論の核物質は等スケール軌跡上で \(J=0\) である。二類がともにその軌跡にいるならば、類の置換は核子の \(J\) を変えない。同じ理論は、\(E^{2}=B^{2}\) の配置を \(J=0\)、\(M>0\) の釣り合った質量として明示し、純電気の配置と区別する。電子の静止は、二つの予算が一致するその配置として読まれている。電荷ラベルは \(J\) の符号とは独立である。したがって電荷 \(-e\) と \(J=0\) は一つのローターの上で両立し、打ち消し用の磁気ローターは釣り合いの条件からは要らない。
\end{proof}

連続スペクトルが要求する沈黙した粒子は、この命題の後にも実験室の事実として残る。命題が除くのは、その粒子を \(J\) の釣り合いから導く読みである。粒子にレプトン数を与えること、それをニュートリノと同定したうえで質量を定めること、および放出のカイラル構造を定めることは、いずれもこの六空間の外にある。親理論は、双対写像がカイラリティ射影ではないこと、および平行な平面にカイラリティを保つワインバーグ角が無いことをすでに示している。本稿はその二つの否定を V--A の代わりには使わない。

自由空間では、置換の向きを質量が決める。自由中性子の静止エネルギーは自由陽子と電子の和より大きい。超過は第\ref{sec:gap}節の間隙から電子の静止を引いたものであり、\(782/1000\,\mathrm{MeV}\) と \(783/1000\,\mathrm{MeV}\) のあいだに厳密にある。したがって自由中性子の置換は開き、自由陽子の逆置換は閉じる。この開閉は質量の入力による。類の代数は、どちらが重いかを選ばない。

核の中では、強い力の同伴稿のクーロン和が向きを変えうる。\(\gamma_E>0\) の各室で、陽子を一つ加えることは既存の陽子対にわたるクーロン和を加え、中性子を加えてもその和は増えない。陽子過剰の積の電磁的な代償が自由な間隙を上回るとき、逆置換は核のエネルギーを下げる。どの核種でその横断が起きるかは、本稿では評価しない。半経験的質量公式の \(Z^{2}/A^{1/3}\) は、強い力の同伴稿がすでに導出しない係数のままである。

\section{質量差の符号}
\label{sec:gap}

CODATA の自由質量は、差 \(m_n-m_p\) を \(12933/10000\,\mathrm{MeV}\) と \(12934/10000\,\mathrm{MeV}\) のあいだに、電子の静止エネルギーを \(51099/100000\,\mathrm{MeV}\) と \(51101/100000\,\mathrm{MeV}\) のあいだに置く~\cite{codata2018}。比 \((m_n-m_p)/m_e\) はしたがって \(5/2\) と \(38/15\) のあいだに厳密にある。下端は、より小さい差をより大きい電子質量で割った \(1.2933/0.51101>5/2\) であり、上端は、より大きい差をより小さい電子質量で割った \(1.2934/0.51099<38/15\) である。区間は比の正体を与えない。

\begin{proposition}
\label{prop:sign}
強い力の同伴稿の長さの範囲では、第一の節における \(ke^{2}/r\) は \(189/100\,\mathrm{MeV}\) より厳密に大きい。自由な質量差はそれより厳密に小さく、符号は逆である。
\end{proposition}

\begin{proof}
同伴稿は、実験室の符号から、第一の節を \(14/19\,\mathrm{fm}\) と \(213/280\,\mathrm{fm}\) のあいだに置く。\(ke^{2}=\alpha\hbar c\) はこの区間の上端で最小になる。\(\alpha\hbar c>1.439\,\mathrm{MeV\,fm}\) であり、
\[
1.439\times\frac{280}{213}>1.891>1.89
\]
である。質量差は \(1.294\,\mathrm{MeV}\) より小さい。正の \(ke^{2}/r\) は荷電側のエネルギーを上げる。自由な順序では荷電の類が軽い。
\end{proof}

したがって類 \([R_{\rm n}]\) の超過は、陽子の電気的自己エネルギーが節に置いた値ではない。自己エネルギーは、観測された順序を打ち消す側に、間隙より大きく働く。超過は類そのものの符号なし質量に属し、クーロン項には属さない。等スケール軌跡の上では、その質量は磁気平方に等しいので、超過はその平方の超過である。絶対値は導出しない。その磁石は第\ref{sec:moment}節である。台帳の \(m_d-m_u\) は、同じ超過をラベルへ割り付けた数であり、ここでも導出されない。

時間尺度も、置換を強いチャンネルから分ける。\(\hbar c<198\,\mathrm{MeV\,fm}\) かつ \(c>2.99\times 10^{23}\,\mathrm{fm\,s^{-1}}\) のとき、\(\hbar<6.63\times 10^{-22}\,\mathrm{MeV\,s}\) である。間隙が \(1.293\,\mathrm{MeV}\) より大きいので、
\[
\frac{\hbar}{(m_n-m_p)c^{2}}<\frac{6.63\times 10^{-22}}{1.293}\,\mathrm{s}<6\times 10^{-22}\,\mathrm{s}
\]
である。報告されている平均寿命は \(870\,\mathrm{s}\) と \(890\,\mathrm{s}\) のあいだにある~\cite{pdg2024}。比は \(10^{24}\) を超える。強い力の同伴稿が扱う MeV の核時間は、この寿命の側には無い。\(\beta\)崩壊は \(\gamma_s\) の節を横切る落下ではない。

\section{電荷に沈黙した磁石}
\label{sec:moment}

等スケール軌跡の上では、磁気平方は符号なし質量に等しく、電気平方もそれに等しい。したがってその軌跡上の正の質量は磁気をもつ。双対角がすべて消える配置は、そこでは質量をもたない。質量の式では、二つの平方がそれぞれ半分を担う。

電磁ラベルはこの等式に入らない。ラベルが零のとき、電荷とベクトルポテンシャルの積を双対角へ足す置き換えは、どの角も動かさない。双対ローターは双対角だけに依存する。ラベルを立てること、あるいは双対角を固定したまま通常セクターの角を変えることは、ローターを書き換えない。類の置換は、ラベルを変え、スピンを残すことができる。

一つの空間軸を保つ生成子は、平行な磁気生成子である。それはローターである。その生成子と平行な電気生成子が張る平面では、軸が固定されたままなのは電気の係数が零のときに限る。平行な電気生成子はローターではない。その一軸に載った双対位相 \(2\pi\) は、ローターを \(-1\) へ送る。双対位相 \(4\pi\) は単位元へ戻す。どちらの核子のスピンも、この半角である。双対角が一致するとき、二つの類はそれを共有する。

静止しているとき、磁気の項は力を及ぼさない。同じ生成子は、物体が動き始めると力を及ぼす。軸に沿った電気生成子は、静止していても力を及ぼす。ラベルが沈黙しているときに残る磁石は、釣り合った質量の双対セクターの平方であり、物体が動き始めてからその物体を押す。これが中性子の微弱な磁石である。

軌跡の上では、中性子類の超過は磁気平方の超過であり、その大きさは自由な質量差に等しい。大きさは入力のままである。それは核磁子の値ではない。

クォークの台帳は、別の座標である。三つのラベルにスピンの重みを割り当てると、台帳は磁気能率の比を決める。標準的な重みでは
\begin{equation}
\label{eq:moments}
\mu_p=\frac43\mu_u-\frac13\mu_d,
\qquad
\mu_n=\frac43\mu_d-\frac13\mu_u
\end{equation}
であり、電荷比と等しい質量から \(\mu_u=-2\mu_d\) と置く。

\begin{lemma}
\label{lem:moments}
\(\mu_u=-2\mu_d\) かつ \(\mu_d\neq 0\) のとき、\eqref{eq:moments} は \(\mu_n/\mu_p=-2/3\) および \((\mu_p+\mu_n)/(\mu_p-\mu_n)=1/5\) を与える。
\end{lemma}

\begin{proof}
代入により \(\mu_p=(3/2)\mu_u\) および \(\mu_n=-\mu_u\) である。比は \(-2/3\) である。和は \(\mu_u+\mu_d=-\mu_d\) であり、差は \((5/3)(\mu_u-\mu_d)=(5/3)(-3\mu_d)=-5\mu_d\) である。和と差の比は \(1/5\) である。
\end{proof}

測定は、この二つの有理数の近くにあり、その上には無い。CODATA の核磁子単位では \(\mu_p\) は \(2.79284\) と \(2.79286\) のあいだに、\(\mu_n\) は \(-1.91305\) と \(-1.91303\) のあいだにある~\cite{codata2018}。比 \(\mu_n/\mu_p\) は、この矩形の上で \(-1.91305/2.79284\) と \(-1.91303/2.79286\) のあいだに厳密にあり、したがって \(-137/200\) と \(-171/250\) のあいだに厳密にある。\(-2/3\) との差は \(1/55\) と \(1/54\) のあいだに厳密にある。和 \(\mu_p+\mu_n\) は \(0.87979\) と \(0.87983\) のあいだに、差 \(\mu_p-\mu_n\) は \(4.70587\) と \(4.70591\) のあいだにあり、その比は \(9/50\) と \(19/100\) のあいだに厳密にある。\(1/5\) からの不足は \(1/80\) と \(1/70\) のあいだに厳密にある。

重み \(4/3\) と \(-1/3\) は \(\mathfrak{so}(2,1)\) の安定化子からは出ず、上の半角でもない。安定化子の唯一のコンパクト生成子は軸を保つ \(B_a\) であり、スピン \(1/2\) の三つの担体を分類しない。本稿は比 \(-2/3\) を改善しない。幾何が与えるのは磁石の担体であり、比は上に記録した台帳の値のままである。

\section{導かれないもの}
\label{sec:open}

次の量は、類の置換からも、台帳の補題からも出ない。

弱い崩壊の寿命と、電子スペクトルの形。フェルミ定数は入力のままである。上の \(10^{24}\) は、寿命が質量間隙の逆時間ではないという比較であり、\(G_F\) の値ではない。

\(g_A/g_V\)。軸方向結合は六空間のサンドイッチからは読めない。サンドイッチが与えるのは \((E,-B)\) のローレンツ四元力であり、親理論はそれを V--A と区別している。

ニュートリノの質量、ヘリシティ、およびレプトン数。命題~\ref{prop:no-third} は、沈黙した粒子を釣り合いから導かない。

三つの世代、カビボ角、および小林--益川行列。三つの軸は世代ではない。順序積は二類しかもたず、フレーバーを変える因子を含まない。

構造関数と比 \(R\)。補題~\ref{lem:chart} は整数電荷の逆算であり、仮想光子が一つの軸に渡す運動量分率ではない。色の個数は軸の個数として読めない。

質量差の絶対値、および \(\alpha_s\)。命題~\ref{prop:sign} は符号と尺度の障害であり、\(1.293\,\mathrm{MeV}\) の導出ではない。等スケール軌跡の上では、間隙は同じく磁気平方の間隙であり、その読みも間隙を評価しない。

核磁子で書いた能率。軌跡は磁気平方を符号なし質量と同一視し、スピンを平行ローターの半角と同一視する。測定された能率と、台帳の値 \(-2/3\) は、その同一視の外に残る。

荷電パイオンの崩壊は、巻き数零の物体が電磁ラベルをレプトンとして出し、核子因子を残さない遷移として同じチャンネルに並ぶ。強い力の同伴稿は、荷電パイオンのコンプトン波長を長さ \(\ell\) の同定としてすでに用いており、その同定は導出ではない。パイオンが電子よりミュー粒子を選ぶことも、ここからは出ない。

\section{結論}
\label{sec:conclusion}

中性子の\(\beta\)崩壊は、順序積の一つの因子を \([R_{\rm n}]\) から \([R_{\rm p}]\) へ置き換える遷移である。因子の個数と巻き数 \(w=3\) は保たれ、その因子の電磁ラベルは \(0\) から \(+e\) へ変わる。電荷を保つ放出は電子である。連続スペクトルが要求する電気的に沈黙した粒子は、六空間に読みの場所をもち、\(J=0\) の釣り合いからは強制されない。自由空間で置換が中性子から陽子へ向かうのは、中性子類が陽子類と電子の和より重いからであり、その不等式は入力である。核の中で逆が開きうるのは、陽子対のクーロン和が自由な間隙を上回りうるからである。

クォーク電荷 \(2/3\) と \(-1/3\) は、二つの核子電荷を三つの加法的ラベルと二つの値で書いたときに残る唯一の有理解である。三つの軸は、そのラベルの三つの粒子ではない。それらは同一の \(3+3\) 分割の代替的な埋め込みであり、色でも世代でもない。巻き数 \(3\) は一つの類が三層で閉じることである。自由な質量差は、第一の節に置いた正のクーロン・エネルギーより小さく、符号は逆である。等スケール軌跡の上では、その間隙は磁気平方の超過である。中性子はその平方を、電磁ラベルを沈黙させたまま担う。したがってベクトルポテンシャルは双対角を動かさず、磁石は静止した物体に力を及ぼさない。スピンは平行な磁気ローターの半角である。双対位相 \(2\pi\) はローターを \(-1\) へ送り、双対位相 \(4\pi\) は単位元へ戻す。双対角を固定するとき、置換はそのローターを残す。クォークの台帳は \(-2/3\) と \(1/5\) を与え、測定はそのどちらよりも小さい側に、それぞれ \(1/55\) と \(1/54\) のあいだ、および \(1/80\) と \(1/70\) のあいだの間隙を残す。これらの比は半角ではない。寿命は質量間隙の逆時間より \(10^{24}\) 以上長い。

弱い相互作用の演算子は、この読みの後にも未導出である。

\begin{thebibliography}{9}

\bibitem{dst-main}
Anonymous,
``Gravity as Torsion between Dual Spacetime: A Biquaternionic Reformulation of General Relativity'' (2025).
\url{https://github.com/hypernumbernet/dual-spacetime-doc/blob/main/double-spacetime-theory.pdf}.

\bibitem{dst-dio}
Anonymous,
``Discrete Biquaternionic Double Spacetime: Rigorous Proofs of Diophantine Conjectures via Torsional Mismatch and PGA Embedding,''
companion note.

\bibitem{dst-strong}
Anonymous,
``The Strong Force as a Repulsive Layer of Proximity Gravity,''
companion note.

\bibitem{dst-shell}
Anonymous,
``Electron Orbits from the Signed Electrostatic Force: Covalent Bonds, Metallic Bonds, and Crystals,''
companion note.

\bibitem{codata2018}
E.~Tiesinga, P.~J.~Mohr, D.~B.~Newell, and B.~N.~Taylor,
``CODATA recommended values of the fundamental physical constants: 2018,''
\textit{Rev.\ Mod.\ Phys.} \textbf{93}, 025010 (2021).

\bibitem{pdg2024}
S.~Navas \textit{et al.} (Particle Data Group),
``Review of Particle Physics,''
\textit{Phys.\ Rev.\ D} \textbf{110}, 030001 (2024).

\end{thebibliography}

\end{document}
